% ============================== CHAPTER 6 ===================================
\chapter{Implementation}
\label{ch:implementation}

This chapter describes how the design of Chapter~\ref{ch:design} was realised in
code. It covers the implementation strategy and the standards applied, presents
the significant source with commentary, and discusses the optimisations that
made the system fit the memory, latency and size budgets. Complete source is not
reproduced --- the repository contains roughly six thousand lines --- but every
component whose behaviour is non-obvious is shown and explained.

\section{Implementation Approaches}
\label{sec:implapproach}

\subsection{Development Strategy}
\label{subsec:strategy}

The project was implemented as a sequence of vertical slices rather than as
horizontal layers. Milestone M0 delivered a complete, working, end-to-end
system in which the translation engine was a passthrough stub: the command-line
tool accepted a language pair, loaded configuration, called an engine and
printed output. Nothing translated anything, but every interface between
components existed and was under test. Each subsequent milestone replaced one
stub with a real implementation, and the test suite grew monotonically.

This ordering was chosen deliberately, and it paid for itself. Because the
inference interface was fixed at M0, the ONNX engine that arrived at M5 was a
drop-in replacement; because \texttt{TranslationResult} was defined at M0, the
pivot router added at Phase~2 only needed one new optional field. The
alternative --- building the training pipeline first and the interfaces last ---
would have deferred every integration problem to the end of the project.

Three standards were applied throughout.

\textbf{Configuration over code.} No language pair, dimension, threshold or
hyperparameter appears as a literal in a module. Everything lives in YAML under
\texttt{configs/}, and a test verifies the pipeline is pair-agnostic by running
it on a pair other than the default.

\textbf{One boundary for the teacher.} The teacher's dependencies may be
imported by exactly one file, enforced by a test that scans the source tree.

\textbf{Test-first for anything subtle.} Mechanisms whose failure is silent were
given tests before implementation. The DPO reference model must stay frozen; the
EWC penalty must be zero at the snapshot and positive away from it; the
scorecard must print \textsc{unverified} rather than a pass when evidence is
absent. Each of these is a property that no amount of manual inspection reliably
catches.

\subsection{Repository Structure}
\label{subsec:repostructure}

\begin{lstlisting}[language=bash, caption={Repository layout}, label={lst:layout}]
SETU/
|-- setu_cli.py              # CLI entry point
|-- configs/                 # languages, model, training, data, teacher, preference
|   |-- model.yaml           #   CPU development student
|   `-- model.gpu.yaml       #   52.6M GPU student
|-- scripts/
|   |-- train_full.py        # teacher ceiling + train + evaluate in one command
|   |-- build_scores.py      # scorecard generation from per-direction reports
|   `-- translation_suite.py # loads each model and checks its output is healthy
|-- src/setu/
|   |-- types.py             # CorpusEntry, PreferencePair, ModelConfig, TranslationResult
|   |-- config.py            # YAML loaders, language registry, ISO <-> FLORES
|   |-- corpus/              # loader, normalize, filters, dedup, pipeline
|   |-- teacher/             # base (the wall) + indictrans2 backend
|   |-- distill/             # SeqKDDistiller + pipeline
|   |-- preference/          # knn, generator, validate, pipeline
|   |-- training/            # tokenizer, student, dataset, sft, dpo, aio_kd, ewc
|   |-- eval/                # metrics, testsets, comet, significance
|   |-- quantize/            # ONNX export + progressive quantisation
|   |-- inference/           # engine, onnx_engine, router
|   `-- report.py            # final scorecard
|-- interfaces/rest/app.py   # FastAPI service
|-- frontend/                # Next.js client and offline PWA
|-- vm/                      # GPU VM runners and training outputs
`-- tests/                   # 85 tests
\end{lstlisting}

\section{Coding Details and Code Efficiency}
\label{sec:coding}

\subsection{The Teacher Abstraction and Its Wall}
\label{subsec:codewall}

Every dependency on IndicTrans2 is confined to one class behind a three-method
interface. The abstraction is narrow on purpose: a distiller needs only to
generate candidates and score them.

\begin{lstlisting}[language=Python, caption={The teacher interface (\texttt{src/setu/teacher/base.py})}, label={lst:teacherbase}]
class TeacherModel:
    """Abstract teacher. Only IndicTrans2Teacher may import transformers/torch."""

    def generate_candidates(self, src_text: str, src_lang: str,
                            tgt_lang: str, n: int = 4) -> list[str]:
        """Return n candidate translations, best first."""
        raise NotImplementedError

    def generate_candidates_batch(self, src_texts: list[str], src_lang: str,
                                  tgt_lang: str, n: int = 4) -> list[list[str]]:
        """Batched form -- the only method the SeqKD distiller needs."""
        raise NotImplementedError

    def score_candidates(self, candidates: list[str],
                         reference: str) -> list[float]:
        """Sentence-level chrF of each candidate against the reference."""
        raise NotImplementedError
\end{lstlisting}

The wall is enforced mechanically. The test below walks the source tree and
fails if any module other than the designated backend imports the teacher's
dependencies. Without it the restriction would erode within weeks, and the
deployed artifact would silently acquire a PyTorch dependency.

\begin{lstlisting}[language=Python, caption={Enforcing the teacher wall (\texttt{tests/test\_teacher.py})}, label={lst:wall}]
FORBIDDEN = ("transformers", "torch", "IndicTransToolkit")
ALLOWED   = {"src/setu/teacher/indictrans2.py"}

def test_teacher_wall():
    """No module outside the backend may import the teacher's dependencies."""
    for path in Path("src/setu").rglob("*.py"):
        rel = path.as_posix()
        if rel in ALLOWED or "/training/" in rel or "/quantize/" in rel:
            continue                      # training and export legitimately use torch
        source = path.read_text()
        for name in FORBIDDEN:
            assert f"import {name}" not in source, f"{rel} imports {name}"
\end{lstlisting}

\subsection{Corpus Cleaning}
\label{subsec:codecorpus}

The cleaning filters encode domain knowledge about Indic text that a generic
text pipeline would get wrong. Two details matter. Zero-width joiners and
non-joiners must survive normalisation, because they distinguish valid
conjuncts; a naive control-character strip corrupts correct text. And the
script-fraction filter is what removes mined pairs in which one side is simply
not in the claimed language.

\begin{lstlisting}[language=Python, caption={Corpus filters (\texttt{src/setu/corpus/filters.py})}, label={lst:filters}]
ZWJ, ZWNJ = "‍", "‌"

def normalize(text: str) -> str:
    """NFC normalise and strip control characters, preserving ZWJ/ZWNJ."""
    text = unicodedata.normalize("NFC", text)
    return "".join(
        ch for ch in text
        if unicodedata.category(ch) != "Cc" or ch in (ZWJ, ZWNJ, "\n", "\t")
    )

def passes_length_ratio(src: str, tgt: str, max_ratio: float = 3.0) -> bool:
    """Reject misaligned pairs: mined corpora pair unrelated sentences."""
    if not src or not tgt:
        return False
    ratio = max(len(src), len(tgt)) / min(len(src), len(tgt))
    return ratio <= max_ratio

def passes_script_fraction(text: str, script: str, minimum: float = 0.5) -> bool:
    """At least `minimum` of the letters must belong to the expected script."""
    letters = [ch for ch in text if ch.isalpha()]
    if not letters:
        return False
    in_script = sum(1 for ch in letters if script_of(ch) == script)
    return in_script / len(letters) >= minimum
\end{lstlisting}

\subsection{The \seqkd{} Distiller}
\label{subsec:codedistill}

The distiller is short, which is appropriate: sequence-level distillation is a
simple idea, and the project's contribution is the evidence that it is the right
one, not the code that implements it.

\begin{lstlisting}[language=Python, caption={Sequence-level distillation (\texttt{src/setu/distill/distiller.py})}, label={lst:distill}]
class SeqKDDistiller:
    """Replace each entry's human target with the teacher's 1-best translation."""

    def __init__(self, teacher: TeacherModel, batch_size: int = 32):
        self.teacher = teacher
        self.batch_size = batch_size
        self.stats = {"entries": 0, "empty": 0}

    def distill(self, entries: list[CorpusEntry]) -> list[CorpusEntry]:
        out: list[CorpusEntry] = []
        for start in range(0, len(entries), self.batch_size):
            chunk = entries[start:start + self.batch_size]
            cands = self.teacher.generate_candidates_batch(
                [e.src_text for e in chunk],
                chunk[0].src_lang, chunk[0].tgt_lang, n=1,
            )
            for e, c in zip(chunk, cands):
                if c and c[0].strip():
                    tgt = c[0]                  # the teacher's target replaces the human one
                else:
                    tgt = e.tgt_text            # safe fallback, never a blank target
                    self.stats["empty"] += 1
                out.append(CorpusEntry(e.src_lang, e.tgt_lang,
                                       e.src_text, tgt, "seqkd"))
            self.stats["entries"] += len(chunk)
        return out
\end{lstlisting}

The \texttt{source} field is set to \texttt{"seqkd"} so that a distilled corpus
can never be mistaken for a reference corpus downstream. The evaluation code
asserts that it is scoring against entries whose source is \emph{not}
\texttt{seqkd}, which is the mechanical guarantee behind the fairness of
Section~\ref{sec:headtohead}.

\subsection{The Training Loop}
\label{subsec:codetrain}

This is the most-debugged code in the project. Every line commented below
corresponds to a failure that was observed and diagnosed; the run history is in
Section~\ref{sec:runlog}.

\begin{lstlisting}[language=Python, caption={Supervised training step with all four stabilisers (\texttt{src/setu/training/sft.py})}, label={lst:sft}]
def train(self, dataset) -> dict:
    total_steps = len(dataloader) * self.epochs
    # Warmup was previously read from config and silently ignored; a constant LR
    # underfits a from-scratch transformer. Cap warmup at 10% of the run.
    warmup = min(self.config.get("warmup_steps", 0), max(1, total_steps // 10))
    scheduler = get_linear_schedule_with_warmup(self.optimizer, warmup, total_steps)

    for epoch in range(self.epochs):
        for batch in dataloader:
            outputs = self.model(input_ids=batch["input_ids"],
                                 attention_mask=batch["attention_mask"],
                                 labels=batch["labels"])   # labels use -100 at PAD
            if self.label_smoothing > 0:
                loss = F.cross_entropy(
                    outputs.logits.view(-1, outputs.logits.size(-1)),
                    batch["labels"].view(-1),
                    ignore_index=-100,                  # never score padding
                    label_smoothing=self.label_smoothing,
                )
            else:
                loss = outputs.loss
            loss.backward()
            # The single most important stabiliser: without clipping, a 200k-sentence
            # run collapsed to constant output while 120k on the same code succeeded.
            torch.nn.utils.clip_grad_norm_(self.model.parameters(), self.max_grad_norm)
            self.optimizer.step()
            scheduler.step()
            self.optimizer.zero_grad()
\end{lstlisting}

The dataset builds the labels, and the single most consequential line in the
entire training path is the padding mask:

\begin{lstlisting}[language=Python, caption={Label construction (\texttt{src/setu/training/dataset.py})}, label={lst:dataset}]
def encode_target(self, text: str) -> list[int]:
    """Labels are [tokens, EOS] with NO leading BOS: HuggingFace already prepends
    the decoder start token, so adding one wastes a decode step and creates a
    train/generate mismatch that produced empty output in an early run."""
    ids = self.sp.encode(text)[: self.max_seq_len - 1]
    return ids + [self.eos_id]

def collate(self, batch: list[dict]) -> dict:
    labels = pad_sequence([b["labels"] for b in batch], batch_first=True,
                          padding_value=PAD_ID)
    # -100, not PAD: HuggingFace's loss uses ignore_index=-100. Padding with PAD
    # gave the model free credit for predicting PAD, producing a fake-low loss of
    # 0.274 and a decoder that ignored the source entirely.
    labels[labels == PAD_ID] = -100
    return {"input_ids": ..., "attention_mask": ..., "labels": labels}
\end{lstlisting}

\subsection{The Inference Router}
\label{subsec:coderouter}

The router is the whole of \setu{}'s multilingual coverage story. It is also the
one place where routing logic exists, which is what prevents the interfaces from
diverging.

\begin{lstlisting}[language=Python, caption={English-pivot routing (\texttt{src/setu/inference/router.py})}, label={lst:router}]
@lru_cache(maxsize=None)
def engine_for_pair(pair: str) -> InferenceEngine:
    """One cached engine per pair; a pair with no model on disk stays a stub."""
    return InferenceEngine(ModelConfig(language_pair=pair))

def translate(text, src_lang, tgt_lang) -> tuple[TranslationResult, str | None, bool]:
    """Returns (result, pivot, stub). Raises ValueError on unknown or equal codes."""
    src, tgt = resolve_language(src_lang), resolve_language(tgt_lang)
    if src["iso"] == tgt["iso"]:
        raise ValueError(f"Source and target language are both {src['iso']!r}")

    direct = f"{src['flores']}-{tgt['flores']}"
    if _available(direct):
        return engine_for_pair(direct).translate(text, src["iso"], tgt["iso"]), None, False

    # No direct model: chain source -> English -> target if both halves exist.
    eng = load_languages()["en"]["flores"]            # eng_Latn
    if src["iso"] != "en" and tgt["iso"] != "en":
        hop1, hop2 = f"{src['flores']}-{eng}", f"{eng}-{tgt['flores']}"
        if _available(hop1) and _available(hop2):
            r1 = engine_for_pair(hop1).translate(text, src["iso"], "en")
            r2 = engine_for_pair(hop2).translate(r1.translated_text, "en", tgt["iso"])
            return (TranslationResult(
                        translated_text=r2.translated_text,
                        src_lang=src["iso"], tgt_lang=tgt["iso"],
                        latency_ms=(r1.latency_ms or 0.0) + (r2.latency_ms or 0.0)),
                    "en", False)
    # Nothing can serve this pair: pass through, but say so.
    return engine_for_pair(direct).translate(text, src["iso"], tgt["iso"]), None, True
\end{lstlisting}

\subsection{Offline Verification}
\label{subsec:codeoffline}

The offline guarantee is a claim about the absence of behaviour, which cannot be
established by inspection alone. It is therefore tested by removing the
capability entirely: all socket construction is patched to raise, and
translation is then required to succeed.

\begin{lstlisting}[language=Python, caption={Proving offline operation (\texttt{tests/test\_quantize.py})}, label={lst:offline}]
def assert_offline():
    """Disable every socket so any network call raises instead of succeeding."""
    def _blocked(*args, **kwargs):
        raise RuntimeError("network access attempted during offline inference")
    socket.socket = _blocked
    socket.create_connection = _blocked

def test_onnx_engine_translates_offline(tmp_models):
    assert_offline()                          # no sockets from here on
    engine = InferenceEngine(ModelConfig(language_pair="hin_Deva-eng_Latn"))
    result = engine.translate("bhaarat ek vishaal desh hai.", "hi", "en")
    assert result.translated_text                       # non-empty
    assert not engine.is_stub                           # a real model answered
\end{lstlisting}

\subsection{The REST Service}
\label{subsec:coderest}

\begin{lstlisting}[language=Python, caption={Translation endpoint (\texttt{interfaces/rest/app.py})}, label={lst:rest}]
@app.post("/translate", response_model=TranslateResponse)
def translate(req: TranslateRequest) -> TranslateResponse:
    try:
        result, pivot, stub = router.translate(req.text, req.source_lang, req.target_lang)
    except ValueError as exc:                 # unknown code, or src == tgt
        raise HTTPException(status_code=400, detail=str(exc))
    return TranslateResponse(
        translated_text=result.translated_text,
        source_lang=req.source_lang, target_lang=req.target_lang,
        latency_ms=result.latency_ms,
        pivot=pivot,      # "en" when routed through English, else null
        stub=stub,        # true means passthrough: never mistake it for a translation
    )

@app.get("/models")
def models() -> dict:
    """Discover every model directory on disk with its variant and size."""
    return {"count": len(_scan_models()), "models": _scan_models()}
\end{lstlisting}

\subsection{Code Efficiency and Optimisation}
\label{subsec:efficiency}

Efficiency in \setu{} means fitting three budgets: training memory, artifact
size and inference latency. Table~\ref{tab:optimisations} records the
optimisations applied and their measured effect.

\begin{longtable}{@{}L{3.4cm} L{5.2cm} L{5.4cm}@{}}
\caption{Optimisations applied and their measured effect}
\label{tab:optimisations}\\
\toprule
\rowcolor{headblue}
\textbf{Optimisation} & \textbf{Problem addressed} & \textbf{Measured effect} \\
\midrule
\endfirsthead
\toprule
\rowcolor{headblue}
\textbf{Optimisation} & \textbf{Problem addressed} & \textbf{Measured effect} \\
\midrule
\endhead
\bottomrule
\endlastfoot

Progressive INT8 then INT4 quantisation & 410.9\,MB FP32 artifact is over twice the 200\,MB budget & 410.9 $\rightarrow$ 104.6 $\rightarrow$ 103.9\,MB, a 4.0$\times$ compression, with negligible quality loss \\
\rowcolor{rowgrey}
Greedy decoding at deployment & Beam search costs roughly 4$\times$ the latency & p90 latency falls to 141--350\,ms for most directions; beam search retained for offline evaluation only \\
Corpus streaming & The full Samanantar corpus is tens of gigabytes & Only the requested slice is ever materialised \\
\rowcolor{rowgrey}
Greedy teacher decoding for distillation & Beam-5 distillation of 500{,}000 sentences does not fit a session & 4--5$\times$ faster; target quality unaffected \\
Per-pair engine cache (\texttt{lru\_cache}) & Reloading a 104\,MB ONNX session per request & Each direction loads from disk at most once \\
\rowcolor{rowgrey}
Free the teacher before training & Teacher plus student exceeded 15\,GB of GPU memory & Eliminated the CUDA out-of-memory failure in the first GPU run \\
Batch size 32 for SFT, 16 for DPO & DPO holds policy and reference and runs four forward passes per step & Both fit a 15--16\,GB accelerator \\
\rowcolor{rowgrey}
16{,}000-token vocabulary & A 32{,}000-token vocabulary left most tokens undertrained & Eliminated degenerate decoding; smaller embedding matrix \\
Sequence truncation at \texttt{max\_seq\_len} & Sentences beyond the positional-embedding limit crashed training and evaluation & Fixed the crash; a regression test prevents recurrence \\
\rowcolor{rowgrey}
Resumable stage markers & A dropped SSH session discarded hours of work & Interrupted runs resume at the incomplete stage \\
\texttt{expandable\_segments} disabled & The CUDA virtual-memory allocator is unsupported on a vGPU slice & Removed the immediate crash on first allocation \\

\end{longtable}

\subsubsection{The Quantisation Trade-off, Measured}

\begin{table}[H]
\centering
\caption{Effect of quantisation on size, latency and quality (Hindi\arrowto{}English student)}
\label{tab:quanttradeoff}
\begin{tabular}{@{}l r r r r@{}}
\toprule
\rowcolor{headblue}
\textbf{Precision} & \textbf{Size (MB)} & \textbf{p90 latency (ms)} & \textbf{BLEU} & \textbf{Within budget} \\
\midrule
FP32 ONNX & 410.9 & 1196 & --- & No (size and latency) \\
\rowcolor{rowgrey}
INT8      & 104.6 & 636 $\rightarrow$ 335 & 18.3 & Yes \\
INT4      & \textbf{103.9} & \textbf{567 $\rightarrow$ 191} & \textbf{19.3} & Yes \\
\bottomrule
\end{tabular}
\end{table}

The arrows record an important intermediate finding. The first quantisation run
produced latencies above the 500\,ms budget not because quantisation was
ineffective but because the undertrained model of that run decoded 128 junk
tokens for every input. Once the training bugs were fixed and the model emitted
sentences of normal length, the same quantised artifact measured 191\,ms. The
lesson is that a latency measurement on a broken model measures the breakage,
not the runtime.

The quality figures also show that INT4 scored marginally \emph{higher} than
INT8 (19.3 against 18.3 BLEU). This is quantisation variance on a
500-sentence evaluation set rather than a real gain from lower precision, but it
confirms that 4-bit quantisation costs nothing measurable here, which is why the
INT4 variant is the one shipped.
